% Web presentation for the Math I notes.
% Chapter TeX files should keep using semantic environments only.

\IfFontExistsTF{TeX Gyre Termes}
  {\setmainfont{TeX Gyre Termes}}
  {\setmainfont{Times New Roman}}
\IfFontExistsTF{TeX Gyre Heros}
  {\setsansfont{TeX Gyre Heros}[Scale=MatchLowercase]}
  {\setsansfont{Arial}[Scale=MatchLowercase]}
\IfFontExistsTF{TeX Gyre Cursor}
  {\setmonofont{TeX Gyre Cursor}[Scale=MatchLowercase]}
  {\setmonofont{Consolas}[Scale=MatchLowercase]}

\IfFontExistsTF{SimSun}
  {\setCJKmainfont{SimSun}[AutoFakeBold=2.5, AutoFakeSlant=0.18]}
  {\setCJKmainfont{Songti SC}[AutoFakeBold=2.5, AutoFakeSlant=0.18]}
\IfFontExistsTF{Microsoft YaHei}
  {\setCJKsansfont{Microsoft YaHei}[AutoFakeBold=2.0]}
  {\setCJKsansfont{SimHei}[AutoFakeBold=2.0]}

\definecolor{mathOneInk}{HTML}{252522}
\definecolor{mathOneMuted}{HTML}{6F6A60}
\definecolor{mathOneRule}{HTML}{D8D1C3}
\definecolor{mathOneForest}{HTML}{1F3B2D}
\definecolor{mathOneGold}{HTML}{A9874A}
\definecolor{mathOneBurgundy}{HTML}{7A2E35}

\color{mathOneInk}
\linespread{1.12}
\setlength{\parindent}{2em}
\setlength{\parskip}{0.35em}
\numberwithin{equation}{section}

\hypersetup{
  colorlinks=true,
  linkcolor=mathOneForest,
  urlcolor=mathOneForest,
  citecolor=mathOneGold,
  filecolor=mathOneForest,
  pdfborder={0 0 0}
}

\NewDocumentCommand{\problemAnchor}{m}{\hypertarget{problem:#1}{}}
\NewDocumentCommand{\problemRef}{m}{\hyperlink{problem:#1}{#1}}
\NewDocumentCommand{\problemIndexAnchor}{m}{\hypertarget{problem-index:#1}{}}
\NewDocumentCommand{\problemIndexRef}{m}{\hyperlink{problem-index:#1}{索引}}
\NewDocumentCommand{\knowledgeAnchor}{O{} m}{%
  \IfBlankTF{#1}{\hypertarget{knowledge:#2}{}}{\hypertarget{knowledge:#1}{}}%
}
\NewDocumentCommand{\knowledgeRef}{O{} m}{%
  \IfBlankTF{#1}{\hyperlink{knowledge:#2}{#2}}{\hyperlink{knowledge:#1}{#2}}%
}
\NewDocumentCommand{\knowledgeIndexAnchor}{O{} m}{%
  \IfBlankTF{#1}{\hypertarget{knowledge-index:#2}{}}{\hypertarget{knowledge-index:#1}{}}%
}
\NewDocumentCommand{\knowledgeIndexRef}{O{} m}{%
  \IfBlankTF{#1}{\hyperlink{knowledge-index:#2}{#2}}{\hyperlink{knowledge-index:#1}{#2}}%
}

\newtheoremstyle{mathoneplain}
  {0.8em}{0.8em}{\itshape}{}{\sffamily\bfseries\color{mathOneForest}}{.}{0.6em}{}
\newtheoremstyle{mathonedefinition}
  {0.8em}{0.8em}{\normalfont}{}{\sffamily\bfseries\color{mathOneForest}}{.}{0.6em}{}
\theoremstyle{mathoneplain}
\newtheorem{theorem}{定理}[section]
\newtheorem{lemma}{引理}[section]
\newtheorem{corollary}{推论}[section]
\theoremstyle{mathonedefinition}
\newtheorem{definition}{定义}[section]

\NewDocumentCommand{\mathOneWebBlockTitle}{m m}{%
  \BlockClassSingle{math1-block-title #1}{\textbf{#2}}%
  \par\nobreak
}

\NewDocumentEnvironment{problemMeta}{}{%
  \begin{BlockClass}{problem-meta}%
  \begingroup
  \footnotesize\color{mathOneMuted}%
  \setlength{\parindent}{0pt}%
  \setlength{\parskip}{0.18em}%
  \raggedright
}{%
  \par\endgroup
  \end{BlockClass}%
  \ignorespacesafterend
}

\NewDocumentEnvironment{problemBox}{}{%
  \begin{BlockClass}{problem-box}%
  \mathOneWebBlockTitle{problem-title}{题目}%
  \ignorespaces
}{%
  \end{BlockClass}%
  \ignorespacesafterend
}

\NewDocumentEnvironment{solutionBox}{}{%
  \begin{BlockClass}{solution-box}%
  \mathOneWebBlockTitle{solution-title}{题解}%
  \ignorespaces
}{%
  \end{BlockClass}%
  \ignorespacesafterend
}

\NewDocumentEnvironment{knowledgeBox}{}{%
  \begin{BlockClass}{knowledge-box}%
  \mathOneWebBlockTitle{knowledge-title}{知识点}%
  \ignorespaces
}{%
  \end{BlockClass}%
  \ignorespacesafterend
}

\NewDocumentEnvironment{mistakeBox}{}{%
  \begin{BlockClass}{mistake-box}%
  \mathOneWebBlockTitle{mistake-title}{易错点}%
  \ignorespaces
}{%
  \end{BlockClass}%
  \ignorespacesafterend
}
